\section{Métodos de penalización}

La penalización tiene por objetivo aproximar un problema con restricciones por una secuencia de problemas irrestrictos, teniendo en vista la aplicación de técnicas computacionales ya desarrolladas. La idea de convertir un problema de optimización con restricciones en un problema (o una secuencia de problemas) sin restricciones es natural y es usada en la construcción de varios algoritmos.

En la penalización externa, a la función objetivo se le añade un término cuyo valor es nulo en todos los puntos factibles y es positivo para los no factibles. En general, la solución de un problema así penalizado está fuera del conjunto factible, lo que explica el nombre de penalización ``externa''. Para obtener convergencia a la solución del problema original (que, obviamente, tiene que ser factible), el parámetro escalar que multiplica el término penalizador tiene que crecer a más infinito. Una excepción importante a esta regla es la \textit{penalización exacta}, donde un valor del parámetro penalizador suficientemente grande (pero fijo) permite la recuperación exacta de una solución del problema original. Típicamente, esto solo es posible a costa de la utilización de términos penalizadores no diferenciables. La penalización exacta será presentada un poco más adelante, en § 4.3.2.

En la penalización interna, a la función objetivo se le añade un término cuyo valor tiende a más infinito cuando los puntos factibles se aproximan a la frontera del conjunto factible del problema. Tenemos entonces una ``barrera'' que no permite a las soluciones de problemas así penalizados dejar el interior del conjunto factible. Esto explica el nombre de penalización ``interna'' o, también, de métodos de ``barreras''. Para obtener convergencia a la solución del problema original (que, típicamente, está justamente en la frontera del conjunto factible), el parámetro escalar que multiplica el término penalizador tiene que decrecer a cero.

Las técnicas de penalización son bien simples y universales. Sin embargo, ellas no son muy usadas en forma pura, debido a la necesidad de la utilización de valores del parámetro penalizador muy grandes (en el caso de la penalización externa) o muy pequeños (en el caso de la penalización interna), lo que tiende a provocar inestabilidad numérica. Mismo así, la penalización tiene una cierta importancia histórica, y es útil para el entendimiento de estrategias más sofisticadas. Además de eso, la penalización exacta proporciona una de las técnicas más eficientes para la globalización de métodos de programación cuadrática secuencial, que será desarrollada en § 4.6.

\subsection{Penalización externa}

Consideremos un problema en formato general
\begin{equation}\label{eq:4.31}
    \min f(x) \quad \text{sujeto a} \quad x \in D,
\end{equation}
donde $\inmap{f}{\R^n}{\R}$ y $D \subset \R^n$ es un conjunto cualquiera. Como ya fue anunciado, introduciremos un término que tiene por objetivo penalizar la violación de las restricciones, i.e., una función $\inmap{\psi}{\R^n}{\R_+}$, continua en $\R^n$, tal que
\begin{equation}\label{eq:4.32}
    \psi(x) =
    \begin{cases}
        0 & \forall x \in D, \\
        > 0 & \forall x \in \R^n \setminus D.
    \end{cases}
\end{equation}

Por ejemplo, si
\begin{equation}\label{eq:4.33}
    D = \{ x \in \R^n \mid h(x) = 0, g(x) \le 0 \},
\end{equation}
es conveniente definir
\[
    \inmap{\Psi}{\R^n}{\R^l \times \R^m},
\]
\begin{equation}\label{eq:4.34}
    \Psi(x) = (h(x), (\max\{0, g_1(x)\}, \dots, \max\{0, g_m(x)\})),
\end{equation}
\begin{equation}\label{eq:4.35}
    \psi(x) = \|\Psi(x)\|^p, \quad x \in \R^n,
\end{equation}
donde $p > 0$. En \eqref{eq:4.35}, además de la norma Euclidiana $\|\cdot\| = \|\cdot\|_2$, pueden ser utilizadas también $\|\cdot\|_p$, o $\|\cdot\|_\infty$ (para $p \in (0, 1)$, $\|\cdot\|_p$ no define una norma, pero en este contexto es conveniente utilizar la misma notación). Estas elecciones resultan en
\begin{equation}\label{eq:4.36}
    \psi(x) = \|\Psi(x)\|_p^p = \|h(x)\|_p^p + \sum_{i=1}^m (\max\{g_i(x), 0\})^p,
\end{equation}
o
\begin{equation}\label{eq:4.37}
    \psi(x) = \|\Psi(x)\|_\infty^p = (\max \{ \|h(x)\|_\infty, 0, g_1(x), \dots, g_m(x) \})^p.
\end{equation}

Notamos que para $h$ y $g$ diferenciables, la función $\psi$ definida por \eqref{eq:4.36} es diferenciable cuando $p > 1$, y no es diferenciable en caso contrario. La función $\psi$ dada por \eqref{eq:4.37} no es diferenciable, mismo cuando $h$ y $g$ son diferenciables. Como veremos más adelante, existen algunas diferencias importantes entre propiedades de penalizadores diferenciables y no diferenciables.

Consideremos una familia de problemas irrestrictos
\begin{equation}\label{eq:4.38}
    \min \varphi(x; c), \quad x \in \R^n,
\end{equation}
donde
\begin{equation}\label{eq:4.39}
    \inmap{\varphi(\cdot; c)}{\R^n}{\R}, \quad \varphi(x; c) = f(x) + c\psi(x),
\end{equation}
y $c > 0$ es el parámetro penalizador.

La idea es que en la medida en que el valor de $c$ aumenta, la violación de las restricciones se vuelve cada vez más cara. Podemos esperar, entonces, que las soluciones de los problemas irrestrictos \eqref{eq:4.38} se aproximen al conjunto factible cuando $c \to +\infty$ y sus puntos de acumulación resuelvan el problema original \eqref{eq:4.31}.

\begin{algorithm}[Método de penalización externa]
    Elegir una función $\inmap{\psi}{\R^n}{\R_+}$, penalizador externo para el conjunto $D$. Elegir $c_0 > 0$. Tomar $k := 0$.
    \begin{enumerate}
        \item Calcular $x^k = x(c_k)$, una solución del problema \eqref{eq:4.38}, donde la función objetivo es dada por la fórmula \eqref{eq:4.39} con $c = c_k$.
        \item Elegir $c_{k+1} > c_k$, tomar $k := k + 1$. Retornar al Paso 1.
    \end{enumerate}
\end{algorithm}

Mencionamos que, por razones obvias, es natural utilizar el punto $x^{k-1}$ como punto inicial del método aplicado para minimizar la función $\varphi(\cdot; c_k)$ en la iteración $k$.

Notemos que la mayoría de los resultados de convergencia de métodos de este tipo se refiere a la secuencia de soluciones globales de los problemas \eqref{eq:4.38}, suponiendo que ellas existan. Por lo tanto, fuera del caso convexo, la aplicación de métodos de penalización puros no viene con convergencia garantizada, ya que en el caso no convexo no hay cómo calcular soluciones globales. El asunto de la utilización de soluciones locales será considerado también, más adelante.

A continuación, veremos algunos ejemplos de aplicación del Algoritmo 4.3.1.

\begin{example}\label{ex:4.3.1}
    Consideremos el problema
    \[
        \min x_1^2 - x_2 \quad \text{sujeto a} \quad x_1 + x_2 - 1 = 0, \ x_1 \ge 0.
    \]
    Este es un problema de programación convexa, y la única solución es $\bar{x} = (0, 1)$.
    
    Eligiendo el término penalizador \eqref{eq:4.36} con $p = 2$, obtenemos
    \[
        \varphi(x; c) = x_1^2 - x_2 + c(x_1 + x_2 - 1)^2 + c(\max\{-x_1, 0\})^2.
    \]
    Como es fácil ver, esta función es convexa, diferenciable, y su (único) minimizador (global) irrestricto es caracterizado por el sistema de ecuaciones
    \[
        0 = \varphi'(x; c) = 
        \begin{pmatrix}
            2x_1 + 2c(x_1 + x_2 - 1) - 2c\max\{-x_1, 0\} \\
            -1 + 2c(x_1 + x_2 - 1)
        \end{pmatrix},
    \]
    donde $\varphi'(x; c)$ es el gradiente de $\varphi(\cdot; c)$ en $x$ (creemos que no debe haber confusión con la notación de la derivada direccional). De donde obtenemos que
    \[
        2x_1 + 1 - 2c\max\{-x_1, 0\} = 0.
    \]
    La última ecuación no posee solución si $x_1 \ge 0$. Resolviendo el sistema para $x_1 < 0$, obtenemos
    \[
        x_1(c) = -1/(2(1+c)), \quad x_2(c) = 1 + 1/(2c) + 1/(2(1+c)).
    \]
    En particular, tenemos que
    \[
        x(c) = (x_1(c), x_2(c)) \to (0, 1) = \bar{x} \quad (c \to +\infty).
    \]
    Observemos que para todo $c > 0$, el punto $x(c)$ no es factible.
    La Figura 4.3.4 ilustra la trayectoria $x(c)$ de las soluciones de los problemas penalizados.
    
    El ejemplo de arriba muestra que, en general, las soluciones de problemas penalizados no son puntos factibles del problema original y, por lo tanto, para métodos de este tipo solo podemos esperar convergencia asintótica. El ejemplo siguiente muestra que, en principio, es posible que alguna solución del problema penalizado sea factible y que, en este caso, ella es la solución del problema original. Resaltamos, no obstante, que tal situación es bastante atípica en el caso de la utilización de penalizadores diferenciables.
\end{example}

\begin{example}\label{ex:4.3.2}
    Consideremos el problema
    \[
        \min x^2 \quad \text{sujeto a} \quad x \ge -1.
    \]
    La solución de este problema es $\bar{x} = 0$.
    
    Eligiendo el término penalizador \eqref{eq:4.36} con $p = 2$, obtenemos
    \[
        \varphi(x; c) = x^2 + c(\max\{-x - 1, 0\})^2.
    \]
    Como es fácil ver, esta función es convexa, diferenciable, y su (único) minimizador (global) irrestricto es caracterizado por
    \[
        0 = \varphi'(x; c) = 2x - 2c\max\{-x - 1, 0\}.
    \]
    La única solución de esta ecuación es $x(c) = 0$.
    Observemos que $x(c)$ es un punto factible y que $x(c) = \bar{x}$.
\end{example}

A continuación, probaremos las propiedades de convergencia del Algoritmo 4.3.1, ya observadas en los Ejemplos 4.3.1 y 4.3.2.

\begin{proposition}\label{prop:4.3.1}
    Sea $\{x^k\}$ una secuencia generada por el Algoritmo 4.3.1, donde $x^k$ es una solución global de \eqref{eq:4.38} con $c = c_k$.
    Entonces, para todo $k$, se tiene que
    \begin{align}
        \varphi(x^{k+1}; c_{k+1}) &\ge \varphi(x^k; c_k), \label{eq:4.40} \\
        \psi(x^{k+1}) &\le \psi(x^k), \label{eq:4.41} \\
        f(x^{k+1}) &\ge f(x^k). \label{eq:4.42}
    \end{align}
\end{proposition}

\begin{proof}
    Por la optimalidad de $x^k$ en \eqref{eq:4.38} con $c = c_k$, tenemos que
    \begin{align}
        \varphi(x^k; c_k) &\le \varphi(x^{k+1}; c_k) \nonumber \\
        &= f(x^{k+1}) + c_k\psi(x^{k+1}) \nonumber \\
        &\le f(x^{k+1}) + c_{k+1}\psi(x^{k+1}) \nonumber \\
        &= \varphi(x^{k+1}; c_{k+1}) \nonumber \\
        &\le \varphi(x^k; c_{k+1}), \label{eq:4.43}
    \end{align}
    donde también utilizamos que $c_{k+1} > c_k$ y $\psi(x^{k+1}) \ge 0$. Acabamos de mostrar \eqref{eq:4.40}.
    
    De \eqref{eq:4.43}, tenemos
    \begin{align*}
        f(x^k) + c_k\psi(x^k) &\le f(x^{k+1}) + c_k\psi(x^{k+1}), \\
        f(x^{k+1}) + c_{k+1}\psi(x^{k+1}) &\le f(x^k) + c_{k+1}\psi(x^k).
    \end{align*}
    Sumando las dos últimas desigualdades, obtenemos
    \[
        (c_k - c_{k+1})\psi(x^k) \le (c_k - c_{k+1})\psi(x^{k+1}),
    \]
    de donde, teniendo en cuenta que $c_{k+1} > c_k$, se sigue \eqref{eq:4.41}.
    
    Ahora, utilizando \eqref{eq:4.41} y \eqref{eq:4.43}, de
    \begin{align*}
        f(x^k) + c_k\psi(x^k) &\le f(x^{k+1}) + c_k\psi(x^{k+1}) \\
        &\le f(x^{k+1}) + c_k\psi(x^k),
    \end{align*}
    obtenemos \eqref{eq:4.42}.
\end{proof}

El resultado siguiente prueba la convergencia asintótica del Algoritmo 4.3.1, además de mostrar que, en general, los puntos $x^k$ generados no son factibles. Con efecto, si el método genera un punto factible, él ya es una solución del problema original. No obstante, es claro que tal situación no puede ser esperada en general.

\begin{theorem}[Convergencia global del método de penalización externa]\label{th:4.3.1}
    Sea $\{x^k\}$ una secuencia generada por el Algoritmo 4.3.1, donde $x^k$ es una solución global de \eqref{eq:4.38} con $c = c_k$.
    Entonces $x^k \in D$ si, y solamente si, $x^k$ es una solución global del problema \eqref{eq:4.31}.
    
    Sean $\inmap{f}{\R^n}{\R}$ y $\inmap{\psi}{\R^n}{\R_+}$ funciones continuas en $\R^n$, y sea $\bar{v} > -\infty$ el valor óptimo del problema \eqref{eq:4.31}. Si $c_k \to +\infty \ (k \to \infty)$, entonces todo punto de acumulación de la secuencia $\{x^k\}$ es una solución global del problema \eqref{eq:4.31}.
\end{theorem}

\begin{proof}
    Primero, es fácil ver que para todo punto $x \in D$,
    \begin{align*}
        f(x^k) &\le f(x^k) + c_k\psi(x^k) \\
        &= \varphi(x^k; c_k) \\
        &\le \varphi(x; c_k) \\
        &= f(x),
    \end{align*}
    donde utilizamos que $c_k > 0$, $\psi(x^k) \ge 0$, que $x^k$ es una solución global de \eqref{eq:4.38} con $c = c_k$, y que $\psi(x) = 0$ para todo $x \in D$. Concluimos que, para todo $k$,
    \begin{equation}\label{eq:4.44}
        f(x^k) \le \varphi(x^k; c_k) \le \bar{v} = \inf_{x \in D} f(x).
    \end{equation}
    Si $x^k \in D$, obtenemos entonces que $x^k$ es una solución del problema \eqref{eq:4.31}.
    
    Por la Proposición 4.3.1, la secuencia $\{\varphi(x^k; c_k)\}$ es no decreciente. Por \eqref{eq:4.44}, ella es limitada superiormente por $\bar{v}$. Luego, ella converge y se tiene que
    \begin{equation}\label{eq:4.45}
        \lim_{k \to \infty} \varphi(x^k; c_k) = \lim_{k \to \infty} (f(x^k) + c_k\psi(x^k)) \le \bar{v}.
    \end{equation}
    
    Sea $\bar{x}$ un punto de acumulación de la secuencia $\{x^k\}$, $\{x^{k_j}\} \to \bar{x} \ (j \to \infty)$. Por la continuidad de $f$, tenemos que $f(x^{k_j}) \to f(\bar{x}) \ (j \to \infty)$. De \eqref{eq:4.45}, obtenemos entonces que
    \[
        \lim_{j \to \infty} c_{k_j}\psi(x^{k_j}) \le \bar{v} - f(\bar{x}).
    \]
    De donde, como $c_k \to +\infty$ y $\psi(x^{k_j}) \ge 0$, se sigue que
    \[
        \lim_{j \to \infty} \psi(x^{k_j}) = 0.
    \]
    Por la continuidad de $\psi$, concluimos que $\psi(\bar{x}) = 0$, i.e., $\bar{x} \in D$.
    
    Ahora, pasando al límite en \eqref{eq:4.44}, obtenemos
    \[
        \lim_{j \to \infty} f(x^{k_j}) \le \bar{v},
    \]
    lo que (teniendo en cuenta que $\bar{x} \in D$), muestra que $\bar{x}$ es una solución de \eqref{eq:4.31}.
\end{proof}

El Teorema 4.3.1 puede ser ligeramente generalizado en relación a las hipótesis de continuidad de las funciones $f$ y $\psi$.